% frontmatter/polish/como-usar.tex -- "Jak korzystać z tego zeszytu": las
% instrucciones para el adulto de frontmatter/como-usar.tex, en polaco.
\section*{Jak korzystać z tego zeszytu}

Ten zeszyt jest dla dziecka w wieku siedmiu lub ośmiu lat, które czyta
już samodzielnie i dodaje i odejmuje w zakresie 20. Nie uczy rachunków
dla dorosłych, tylko codziennej ekonomii: że rzeczy mają wartość, że
pieniądze pozwalają wymieniać jedne rzeczy na inne, że nie można mieć
wszystkiego i trzeba wybierać, że oszczędzać to odkładać coś dziś na
jutro, że ktoś piecze chleb, który kupujemy, i że za chodniki i parki
płacimy wszyscy razem. Należy do tej samej rodziny co hiszpańskie
zeszyty \emph{Aprendo a leer} i \emph{Aprendo los números}, z tymi
samymi bohaterami (Lucía, jej młodszy brat Dani, pies Toby, babcia Rosa,
Sofía i Hugo). A strona po stronie jest to polskie wydanie
\emph{Aprendo economía}: ta sama historyjka, te same rysunki i to samo
zadanie każdego dnia, więc dziecko może korzystać z jednego zeszytu albo
z obu.

Na każdy dzień roboczy w roku, od poniedziałku do piątku, także w
wakacje, jest jedna strona. Dziesięć minut dziennie w zupełności
wystarczy. Każda strona wygląda tak samo:

\begin{itemize}[leftmargin=6mm]
\item \textbf{Data}, \textbf{Zrobione} i \textbf{Uwagi}: dla dorosłego.
  Zaznaczenie codziennie, czy zadanie było zrobione samodzielnie, z
  pomocą czy z trudem, zajmuje chwilę, a po roku jest zapisem
  wszystkiego, czego dziecko się nauczyło.
\item \textbf{Dzisiaj}, w zielonej ramce: historyjka dnia, dwa albo trzy
  zdania, które czyta dziecko. W poniedziałki pod spodem jest
  \textbf{słowo tygodnia} i jego znaczenie; w piątek się je powtarza, a
  na końcu zeszytu wszystkie są zebrane w \emph{Moim słowniczku
  ekonomicznym}.
\item \textbf{Zadanie}. Polecenia (małą, szarą czcionką) może przeczytać
  dorosły, jeśli trzeba.
\end{itemize}

Każde zadanie jest opisane, jedno po drugim, na następnej stronie.

\subsection*{Jak pomagać}

{\small
\begin{itemize}[leftmargin=6mm, itemsep=1pt, topsep=2pt]
\item Rozmawiajcie o historyjce dnia: co zrobilibyście na miejscu
  bohaterów? W ekonomii prawie nigdy nie ma jednej odpowiedzi; dziecko
  uczy się myśleć o tym, co zyskuje, a z czego rezygnuje.
\item Rachunki z pieniędzmi są łatwiejsze z prawdziwymi monetami:
  policzcie te, które macie w domu, bawcie się w sklep, a na zakupach
  pozwólcie dziecku zapłacić i sprawdzić resztę.
\item Pomyłka przy małych pieniądzach to najlepszy sposób, żeby nauczyć
  się dobrze z nich korzystać, kiedy będzie ich więcej.
\item Rodzina Lucíi mieszka w Hiszpanii, więc pieniądze w zeszycie to
  euro, zawsze całe: 5 € to pięć euro. Jeśli w domu płacicie złotymi, to
  dobra okazja, żeby porozmawiać o tym, że różne kraje mają różne
  pieniądze.
\item Na koniec każdej części jest medal, a na końcu zeszytu dyplom i
  \textbf{odpowiedzi} do stron, które je mają.
\end{itemize}}
\newpage
